\chapter{Introduction to Exception Handling in Python (Revisited)}

\section{Introduction: Real-Life Hook}

Rafi is a high school student who created a small calculator program.
One day, his friend entered \texttt{0} as the divisor.
Suddenly, the program crashed and showed a scary error message.

Rafi felt confused and worried.
``Why did my program stop? Why didn’t it handle the mistake?''

In real life:
\begin{itemize}
    \item Students type wrong input
    \item Users make mistakes
    \item Files may not exist
\end{itemize}

A good program should not crash.
It should \textbf{handle errors politely}.
This is where \textbf{exception handling} comes in.

\section{Concept: What Is an Exception?}

\begin{definition}{Exception}
An exception is an \textbf{error that occurs during program execution} and stops the normal flow of the program.
\end{definition}

Common examples:
\begin{itemize}
    \item Dividing by zero
    \item Entering text instead of a number
    \item Opening a file that does not exist
\end{itemize}

Example without handling:

\begin{verbatim}
x = 10
y = 0
print(x / y)
\end{verbatim}

This causes a \textbf{runtime error} and stops the program.

\section{Why Exception Handling Is Important}

\begin{itemize}
    \item Prevents program crashes
    \item Improves user experience
    \item Makes programs reliable and professional
    \item Helps in real-world applications like banking, apps, and websites
\end{itemize}

\section{The try and except Block}

Python provides \texttt{try} and \texttt{except} to handle errors.

\begin{verbatim}
try:
    x = int(input("Enter a number: "))
    y = int(input("Enter another number: "))
    print(x / y)
except:
    print("Something went wrong!")
\end{verbatim}

How it works:
\begin{itemize}
    \item Python tries the code inside \texttt{try}
    \item If an error occurs, \texttt{except} runs
    \item Program does not crash
\end{itemize}

\section{Handling Specific Exceptions}

It is better to handle \textbf{specific errors}.

\begin{verbatim}
try:
    num = int(input("Enter a number: "))
    print(10 / num)
except ZeroDivisionError:
    print("You cannot divide by zero.")
except ValueError:
    print("Please enter a valid number.")
\end{verbatim}

Common exceptions:
\begin{itemize}
    \item \texttt{ZeroDivisionError}
    \item \texttt{ValueError}
    \item \texttt{FileNotFoundError}
\end{itemize}

\section{Using else and finally}

\subsection{else Block}

\begin{verbatim}
try:
    num = int(input("Enter a number: "))
    print(10 / num)
except ZeroDivisionError:
    print("Division by zero is not allowed.")
else:
    print("Calculation successful!")
\end{verbatim}

\subsection{finally Block}

\begin{verbatim}
try:
    file = open("data.txt", "r")
    print(file.read())
except FileNotFoundError:
    print("File not found.")
finally:
    print("Program finished.")
\end{verbatim}

\begin{itemize}
    \item \texttt{else} runs when no error occurs
    \item \texttt{finally} always runs
\end{itemize}

\section{Mini-Project: Safe Calculator}

\textbf{Problem:}  
Create a calculator that never crashes, even if the user enters wrong input.

\begin{verbatim}
try:
    a = int(input("Enter first number: "))
    b = int(input("Enter second number: "))
    print("Result:", a / b)
except ZeroDivisionError:
    print("Error: Division by zero.")
except ValueError:
    print("Error: Invalid input.")
else:
    print("Thank you for using the calculator.")
\end{verbatim}

This project teaches:
\begin{itemize}
    \item User input validation
    \item Error handling
    \item Professional program behavior
\end{itemize}

\section{Diagram: Exception Handling Flow}

\begin{center}
\begin{tikzpicture}[node distance=1.5cm]
\node (start) [draw, ellipse, fill=green!20] {Start Program};
\node (try) [draw, rectangle, fill=yellow!20, below of=start] {Try Block};
\node (error) [draw, diamond, fill=orange!20, below of=try] {Error Occurs?};
\node (except) [draw, rectangle, fill=red!20, right of=error, xshift=3cm] {Except Block};
\node (else) [draw, rectangle, fill=green!20, left of=error, xshift=-3cm] {Else Block};
\node (finally) [draw, rectangle, fill=blue!20, below of=error, yshift=-1.5cm] {Finally Block};
\node (end) [draw, ellipse, fill=green!20, below of=finally] {End Program};

\draw[->] (start) -- (try);
\draw[->] (try) -- (error);
\draw[->] (error) -- node[above]{Yes} (except);
\draw[->] (error) -- node[above]{No} (else);
\draw[->] (except) |- (finally);
\draw[->] (else) |- (finally);
\draw[->] (finally) -- (end);
\end{tikzpicture}
\end{center}

\section{Key Concepts}

\begin{itemize}
    \item Exceptions are runtime errors
    \item \texttt{try} tests risky code
    \item \texttt{except} handles errors
    \item \texttt{else} runs when no error occurs
    \item \texttt{finally} always executes
\end{itemize}

\section{End-of-Chapter Exercises}

\subsection{Very Short Questions}

\begin{enumerate}
    \item What is an exception?
    \item Name one built-in Python exception.
    \item Which block always executes?
\end{enumerate}

\subsection{Short Answer Questions}

\begin{enumerate}
    \item Explain the difference between \texttt{try} and \texttt{except}.
    \item Write a program to handle division by zero.
    \item Why is exception handling important in real-life programs?
\end{enumerate}

\subsection{Activity-Based Exercises}

\begin{enumerate}
    \item Create a program that asks for age and handles non-numeric input.
    \item Write a file-reading program that handles missing files.
    \item Modify the calculator to handle multiplication and subtraction errors.
\end{enumerate}

\section{Teacher Notes}

\marginnote{
\textbf{Learning Objectives}
\begin{itemize}
    \item Students understand runtime errors
    \item Students can use try-except blocks
    \item Students can write safer programs
\end{itemize}

\textbf{Teaching Tips}
\begin{itemize}
    \item Demonstrate common mistakes live in class
    \item Encourage students to break programs and fix them
    \item Use real-life examples of user mistakes
\end{itemize}

\textbf{Assessment}
\begin{itemize}
    \item Correct use of exception handling
    \item Clear error messages
\end{itemize}
}